\scanpage{110}
% archive sheet number: 109
donc, pour cette raison, savoir définir les $\mathbb R f_*$, $\mathbb R_! f_*$ par des arguments qui ne sont pas \unsure{des|du|dans} $D^+$, \inserted{comme foncteurs $D^-(X)$ --- $D^-(Y)$, $D(X)$ --- $D(Y)$.} Pour $\mathbb R f_*$, ce sera possible dès que $f_*$ est de dim cohomologique finie. C'est peut-être vrai dès que $Y$ est noethérien, cela \deleted{\unsure{est}} \inserted{résulte} en tous cas de la \unsure{finitude} de la dim ch. si $Y$ est, \deleted{soit de} soit de type fini sur un corps, soit de type fini sur $\operatorname{Spec}\mathbb Z$. Pour $\mathbb R_! f_*$, on est en meilleure posture encore, grâce au changt de base propre, donc \uline{si} $Y$ \uline{qu. cpct}. on a
\diagram{\[
\left.
\begin{aligned}
\mathbb R_! f_*:\quad&D(X)\;\text{---}\;D(Y)\\
&D^+(X)\;\text{---}\;D^+(Y)\\
&D^-(X)\;\text{---}\;D^-(Y)\\
&D^b(X)\;\text{---}\;D^b(Y)
\end{aligned}
\right\}\quad\begin{gathered}\text{si }Y\text{ qu. cpct}\\ f\text{ compactifiable}\end{gathered}
\]}

Comme on y est, notons
\diagram{\[
\left.
\begin{aligned}
\mathbb R f_*:\quad&D(X)\;\text{---}\;D(Y)\\
&D^{\unsure{+|e}}(X)\;\text{---}\;D^{\unsure{+|c}}(Y)\\
&D^-(X)\;\text{---}\;D^-(Y)\\
&D^b(X)\;\text{---}\;D^b(Y)
\end{aligned}
\right\}\quad\begin{gathered}\text{si }Y\text{ de t.f.}\\\text{sur un corps}\\\text{ou sur }\operatorname{Spec}\mathbb Z,\\\text{et }f\text{ de t.f.}\end{gathered}
\]}
\diagram{\[
\begin{aligned}
\mathbb L f^*:\quad&D(\deleted{\unsure{X}}Y)\;\text{---}\;D(X)\\
&D^-(Y)\;\text{---}\;D^-(X)\\
&D^+(Y)\;\text{---}\;D^+(X)\qquad D^b(Y)\to D^b(X)
\end{aligned}
\]}
